\section{Introduction}
\label{sec:intro}

The widespread use of machine learning to make consequential decisions
about individual citizens (including in domains such as credit,
employment, education and criminal sentencing~\cite{hiring, lending,
  policing, sentencing}) has been accompanied by increased reports of
instances in which the algorithms and models employed can be unfair or
discriminatory in a variety of ways~\cite{propublica,Sweeney13}.  As a
result, research on fairness in machine learning and statistics has
seen rapid growth in recent years ~\cite{PedreshiRT08, KamiranCP10,
  CaldersV10, KamiranC11, LuongRT11, KamiranKZ12, KamishimaAAS12,
  HajianD13, FeldmanFMSV15,FishKL16,AdlerFFRSSV16, HardtPS16,
  Chouldechova17, CaldersKKAZ13,DworkHPRZ12}, and several mathematical
formulations have been proposed as metrics of (un)fairness for a
number of different learning frameworks.  While much of the attention
to date has focused on (binary) classification settings, where
standard fairness notions include equal false positive or negative
rates across different populations, less attention has been paid to
fairness in (linear and logistic) regression settings, where the
target and/or predicted values are continuous, and the same value may
not occur even twice in the training data.

In this work, we introduce a rich family of fairness metrics for
regression models that take the form of a fairness regularizer and
apply them to the standard loss functions for linear and logistic
regression.  Since these loss functions and our fairness regularizer
are convex, the combined objective functions obtained from our
framework are also convex, and thus permit efficient
optimization. Furthermore, our family of fairness metrics covers the
spectrum from the type of {\em group\/} fairness that is common in
classification formulations (where e.g. false arrests in one racial
group can be ``compensated'' for by false arrests in another racial
group) to much stronger notions of {\em individual\/} fairness (where
such cancellations are forbidden, and every injustice is charged to
the model). Intermediate fairness notions are also covered. Our
framework also permits one to either forbid the use of a ``protected''
variable (such as race), by demanding that a single model be learned
across all groups, or to build different group-dependent models.

Most importantly, by varying the weight on the fairness regularizer,
our framework permits us to compute the entire ``Pareto curve'' or
efficient frontier of the trade-off between predictive accuracy and
fairness. Such curves are especially important to examine and
understand in a domain-specific manner: since demanding fairness of
models will always come at a cost of reduced predictive
accuracy~\cite{Corbett-DaviesP17, FishKL16, ZafarVGG17, jelveh2015},
it behooves practitioners working with fairness-sensitive data sets to
understand just how mild or severe this trade-off is in their
particular arena, permitting them to make informed modeling and policy
decisions.

Our central results take the form of an extensive comparative
empirical case study across six distinct datasets in which fairness is
a primary concern. For each of these datasets, we compute and examine
the corresponding fairness-accuracy efficiency frontier. We introduce
an intuitive quantity called the {\em Price of Fairness (PoF)\/},
which numerically quantifies the extent to which increased fairness
degrades accuracy. We compare the PoF across datasets, fairness
notions, and treatments of protected variables.

Our primary contributions are:
\begin{itemize}
\item The introduction of a flexible but convex family of fairness
  regularizers of varying strength that spans the spectrum from group
  to individual fairness.
\item The introduction of a quantitative, data-dependent measure of
  the severity of the accuracy-fairness tradeoff.
\item An extensive empirical comparative study across six
  fairness-sensitive data sets.
\end{itemize}

While our empirical study does reveal some reasonably consistent
findings across datasets (e.g. efficiency curves show broadly similar
shapes; PoF generally higher for individual fairness than group;
somewhat surprisingly, PoF not generally improved much when using
protected variables), perhaps the most important message is a
cautionary one: the detailed trade-off between accuracy and fairness,
and the comparison of different fairness notions, appears to be quite
domain-dependent and lacking prescriptive ``universals''. This is
perhaps consistent with the emerging theoretical literature
demonstrating the lack of a single ``right'' definition of
fairness~\cite{KleinbergMR16,Chouldechova17,FriedlerSV16}, and our
work adds evidence to the view that fairness is a topic demanding
careful domain-specific considerations.
